\section{Εσωτερική \ENG{Layered} Αρχιτεκτονική \ENG{Runtime Core}}
\label{sec:layered_runtime_core}

Στο ανώτερο επίπεδο βρίσκονται τα \ENG{API} στρώματα, δηλαδή οι \ENG{controllers} που εκθέτουν το \ENG{REST} συμβόλαιο. Οι \ENG{controllers} είναι \ENG{"thin"}: επικυρώνουν το αίτημα, ανακτούν τυχόν \ENG{token payloads} από το \ENG{security filter} και μεταβιβάζουν την εργασία σε \ENG{orchestrators} ή \ENG{command handlers}. Το δεύτερο επίπεδο είναι το \ENG{orchestration layer}, όπου ανήκουν κλάσεις όπως οι \texttt{\ENG{LaunchOrchestrator}}, \texttt{\ENG{RuntimeOrchestrator}}, \texttt{\ENG{CreateLaunchCommandHandler}} και \texttt{\ENG{TerminateLaunchCommandHandler}}. Εκεί βρίσκονται οι διασταυρώσεις ανάμεσα σε \ENG{domain} κανόνες, προσωρινή κατάσταση, έλεγχο ακολουθιών και μόνιμη αποθήκευση.

Κάτω από το \ENG{orchestration layer} βρίσκεται το στρώμα υπηρεσιών ανάγνωσης και εγγραφής που συνδέεται με το \ENG{domain} μοντέλο και τη μόνιμη αποθήκευση. Οι \texttt{\ENG{CourseReadService}}, \texttt{\ENG{UserReadService}}, \texttt{\ENG{AttemptReadService}}, \texttt{\ENG{AttemptPersistService}}, \texttt{\ENG{LaunchPersistService}} και οι αντίστοιχες \ENG{CRUD} οντότητες υλοποιούν το κύριο σχήμα δεδομένων του συστήματος. Από το \ENG{migration} \texttt{\ENG{V1\_\_init\_schema.sql}} καθιστά σαφές ότι τα \ENG{attempts} αποθηκεύουν την τελική κατάσταση κανονικοποιημένης προόδου, ενώ τα \texttt{\ENG{attempt\_runtime\_snapshots}} λειτουργούν ως ιστορικό καταγραφής των συγχρονισμών της κατάστασης εκτέλεσης προς μόνιμη αποθήκευση.

Η ανάλυση του \ENG{import flow} δείχνει επίσης ότι το \ENG{layered} σχήμα εφαρμόζεται και πριν από την εκτέλεση. Ο \texttt{\ENG{ManifestParser}} ανήκει στο \ENG{domain} στρώμα και αναλαμβάνει αποκλειστικά την ανάγνωση του \texttt{\ENG{imsmanifest.xml}}, την ανίχνευση του προτύπου, την επιλογή \ENG{entrypoint} και τη συλλογή στοιχειώδους μεταδεδομένης πληροφορίας. Δεν κατασκευάζει πλήρες \ENG{activity tree}, ούτε εξάγει \ENG{sequencing rules}. Έτσι, το \ENG{import} στάδιο παράγει ένα λειτουργικό αλλά σκόπιμα περιορισμένο \ENG{course model}, το οποίο αρκεί για το \ENG{runtime} που πράγματι υλοποιείται.

Το Σχήμα~\ref{fig:layered_runtime_core} δείχνει τα βασικά εσωτερικά επίπεδα και την κατεύθυνση εξάρτησης ανάμεσά τους. Η πρακτική συνέπεια αυτής της οργάνωσης είναι ότι ο \ENG{engine} μπορεί να μεταβάλει τους μηχανισμούς αποθήκευσης, \ENG{token handling} ή \ENG{adapter logic} χωρίς να αλλοιώνει το εξωτερικό \ENG{REST API} , αρκεί να διατηρούνται σταθερές οι εσωτερικές ευθύνες κάθε στρώματος.

\begin{figure}[H]
  \centering
\begin{chapterfigurebox}
  \begin{tikzpicture}[
      font=\small,
      box/.style={draw, rounded corners, fill=gray!5, align=center, text width=10.8cm, minimum height=1.05cm},
      emph/.style={draw, rounded corners, fill=gray!12, align=center, text width=10.8cm, minimum height=1.05cm, font=\small\bfseries},
      arrow/.style={-Latex, thick}
    ]
    \node[box] (api) at (0,0) {\ENG{API layer: controllers, request validation, token payload extraction}};
    \node[emph] (orch) at (0,-1.6) {\ENG{Orchestration layer: command handlers, launch/runtime flow control, sequencing of services}};
    \node[box] (domain) at (0,-3.2) {\ENG{Domain / interpretation layer: \texttt{\ENG{ManifestParser}}, \ENG{standard adapters}, \ENG{normalized progress logic}}};
    \node[box] (persist) at (0,-4.8) {\ENG{Persistence services: \ENG{read/persist services for courses, users, enrollments, attempts and launches}}};
    \node[box] (stores) at (0,-6.4) {\ENG{Storage layer: \ENG{Postgres}, \ENG{Redis}, \ENG{MinIO} and supporting infrastructure}};

    \draw[arrow] (api) -- (orch);
    \draw[arrow] (orch) -- (domain);
    \draw[arrow] (domain) -- (persist);
    \draw[arrow] (persist) -- (stores);
  \end{tikzpicture}
\end{chapterfigurebox}
  \caption{Εσωτερική \ENG{layered} οργάνωση του \ENG{engine} με διαδοχικά επίπεδα ευθύνης και κατεύθυνση εξαρτήσεων προς τα κατώτερα στρώματα.}
  \label{fig:layered_runtime_core}
\end{figure}

